\section*{第 \ref{ch:b1:linmaps} 章——线性映射}

\begin{solution}{exo:b1:linmaps:1}
（1）线性：\omterm{prop:b1:vspaces:coordinates}{坐标}是线性表达式。（2）非线性：
$u(2(1,1)) = 4 \neq 2 = 2u(1,1)$。（3）线性：求导与
乘以 $X$ 均为线性，\omterm{def:b1:linmaps:def}{线性映射}之和亦然。（4）线性：
求值与逐点运算相容。
\end{solution}

\begin{solution}{exo:b1:linmaps:2}
核：解方程组 $x + y - z = 0$，$2x + y + z = 0$，$3x + 2y = 0$。
由第三式得 $y = -\frac{3x}{2}$；第一式给出 $z = x + y =
-\frac x2$；代入第二式验证：$2x - \frac{3x}{2} - \frac x2 = 0$，
成立。故 $\ker u = \operatorname{Vect}\bigl((2, -3, -1)\bigr)$
（取 $x = 2$），维数为 $1$。

由秩–零化度定理，$\operatorname{rk} u = 3 - 1 = 2$。像由
典范基的像张成：$u(e_1) = (1,2,3)$，$u(e_2) =
(1,1,2)$，$u(e_3) = (-1,1,0)$；前两个向量\omterm{def:b1:vspaces:free}{线性无关}，且秩
为 $2$，基为 $\bigl((1,2,3), (1,1,2)\bigr)$。

不是\omterm{def:b1:logic:inj}{单射}（$\ker \neq \{0\}$），也不是\omterm{def:b1:logic:inj}{满射}（秩 $2 < 3$）：
与\cref{cor:b1:linmaps:samedim}一致。
\end{solution}

\begin{solution}{exo:b1:linmaps:3}
\emph{核：}$P = P'$ 蕴含 $\deg P = \deg P'$，除非 $P = 0$；
但 $\deg P' < \deg P$ 对 $P \neq 0$ 成立，故 $\ker u = \{0\}$，
于是 $u$ 作为\omterm{def:b1:findim:def}{有限维}空间 $\R_n[X]$ 上的单自\omterm{def:b1:structures:morphism}{同态}，是
同构（\cref{cor:b1:linmaps:samedim}）。

\emph{逆：}设 $v(P) = P + P' + P'' + \dots + P^{(n)}$（在 $\R_n[X]$ 上的
有限和）。则
\[
v\bigl(u(P)\bigr) = \sum_{k=0}^{n} (P - P')^{(k)}
= \sum_{k=0}^{n} P^{(k)} - \sum_{k=0}^{n} P^{(k+1)}
= P - P^{(n+1)} = P ,
\]
裂项相消，因 $P^{(n+1)} = 0$。故 $v = u^{-1}$。
\end{solution}

\begin{solution}{exo:b1:linmaps:4}
$\operatorname{im}(v \circ u) = v(\operatorname{im} u) \subseteq
\operatorname{im} v$：秩 $\leq \operatorname{rk} v$。又 $v$
限制在 $\operatorname{im} u$ 上的像为
$\operatorname{im}(vu)$，由秩–零化度定理，在
$\operatorname{im} u$ 中 $\operatorname{rk}(vu) \leq
\dim\operatorname{im} u = \operatorname{rk} u$。
\end{solution}

\begin{solution}{exo:b1:linmaps:5}
$u^2 = 0$ 表示 $u(u(x)) = 0$ 对所有 $x$ 成立：每个 $u(x)$ 都属于
$\ker u$，即 $\operatorname{im} u \subseteq \ker u$。于是由
秩–零化度定理：
\[
\dim E = \dim\ker u + \operatorname{rk} u \geq 2\operatorname{rk} u .
\]
取等情形的例子：在 $\R^2$ 中取 $u(x, y) = (y, 0)$，则 $u^2 = 0$，
$\operatorname{rk} u = 1 = \frac{\dim E}{2}$。
\end{solution}

\begin{solution}{exo:b1:linmaps:6}
$(pq)^2 = pqpq = ppqq = pq$（由交换性），即为\omterm{def:b1:linmaps:projection}{投影}
（\cref{thm:b1:linmaps:projchar}）。

像：$\operatorname{im}(pq) \subseteq \operatorname{im} p$，且
$= \operatorname{im}(qp) \subseteq \operatorname{im} q$：从而包含于
交集之中。反之，若 $x \in \operatorname{im} p \cap
\operatorname{im} q$，则 $p(x) = x$ 且 $q(x) = x$（不动点
刻画\omterm{def:b1:linmaps:projection}{投影}的像），故 $pq(x) = x$，即 $x \in
\operatorname{im}(pq)$。

核：$\ker p \subseteq \ker(qp) = \ker(pq)$，同理 $\ker q
\subseteq \ker(pq)$，故和包含于其中。反之，设 $pq(x) =
0$，并写出
\[
x = \underbrace{q(x)}_{\in\, \ker p} +
\underbrace{(x - q(x))}_{\in\, \ker q} :
\]
第一项满足 $p(q(x)) = 0$，故属于 $\ker p$；
第二项属于 $\ker q$，因 $q(x - q(x)) = q(x) - q^2(x) = 0$。由此
得 $x \in \ker p + \ker q$。
\end{solution}

\begin{solution}{exo:b1:linmaps:7}
首先注意一般的包含关系 $\ker u \subseteq \ker u^2$ 与
$\operatorname{im} u^2 \subseteq \operatorname{im} u$，并由
秩–零化度定理，（2）$\iff$ （3）（核相等 $\iff$ 秩相等
$\iff$ 像相等，依上述包含关系）。

（1 $\Rightarrow$ 2）：设 $u^2(x) = 0$，则 $u(x) \in \ker u \cap
\operatorname{im} u = \{0\}$，故 $x \in \ker u$。

（2 $\Rightarrow$ 1）：由格拉斯曼公式及秩–零化度定理，$\dim(\ker u +
\operatorname{im} u) = \dim\ker u + \operatorname{rk} u -
\dim(\ker u \cap \operatorname{im} u) = \dim E - \dim(\ker u \cap
\operatorname{im} u)$：和为 $E$ 当且仅当交集为
$\{0\}$。设 $y \in \ker u \cap \operatorname{im} u$：$y = u(x)$ 且
$u(y) = 0$，故 $u^2(x) = 0$，由（2）得 $u(x) = 0$：$y = 0$。由此
得 $E = \ker u \oplus \operatorname{im} u$。
\end{solution}

\begin{solution}{exo:b1:linmaps:8}
若 $\varphi = 0$，则 $\ker\varphi = E \subseteq \ker\psi$，蕴含
$\psi = 0 = 0\cdot\varphi$。若 $\varphi \neq 0$，则 $\ker\varphi$ 是
\omterm{def:b1:linmaps:forms}{超平面}；取 $a \notin \ker\varphi$。令 $\lambda =
\frac{\psi(a)}{\varphi(a)}$。\omterm{def:b1:linmaps:forms}{线性型} $\psi - \lambda\varphi$
在 $\ker\varphi$ 上为零（二者皆然，由包含关系），且在 $a$ 处为零：
它在 $\ker\varphi \oplus Ka = E$ 上为零。于是 $\psi =
\lambda\varphi$。
\end{solution}

\begin{solution}{exo:b1:linmaps:9}
假设 $\lambda_0 x + \lambda_1 u(x) + \dots + \lambda_{n-1}
u^{n-1}(x) = 0$。作用 $u^{n-1}$：所有含因子 $u^{\geq
n}$ 的项皆消失，剩下 $\lambda_0 u^{n-1}(x) = 0$，于是 $\lambda_0 = 0$。
对剩余关系式作用 $u^{n-2}$：$\lambda_1 u^{n-1}(x) =
0$，于是 $\lambda_1 = 0$，依此类推。该族\omterm{def:b1:vspaces:free}{线性无关}，且元素个数为
$n = \dim E$，故为基（见\cref{prop:b1:findim:twoofthree}）。（在
该基下，$u$ 表现为移位——即极大幂零性的
模型）。
\end{solution}

\begin{solution}{exo:b1:linmaps:10}
\begin{enumerate}
  \item 由 $f \circ f = -\mathrm{id}$ 知 $f$ 为\omterm{def:b1:logic:inj}{双射}，故
        可逆（改编自\cref{prop:b1:logic:comp}：$f$ 的双侧
        逆为 $-f$）。若 $f(x) = \lambda x$，$x \neq 0$：
        两边作用 $f$ 得 $-x = \lambda^2 x$，故 $\lambda^2 = -1$：
        在 $\R$ 中不可能。
  \item 逐步构造该族。取 $x_1 \neq 0$：$(x_1,
        f(x_1))$ 由 (1) 知\omterm{def:b1:vspaces:free}{线性无关}。记 $\operatorname{Vect}$ 为张成运算，
        则当前族 $\bigl(x_1, f(x_1), \dots, x_k,
        f(x_k)\bigr)$ 的张成记为 $V_k$——一个在 $f$ 下稳定的\omterm{def:b1:vspaces:subspace}{子空间}
        （每个生成元都映为另一生成元或其相反数：
        $f(f(x_i)) = -x_i$）——若其不等于 $E$，取 $x_{k+1}
        \notin V_k$。\emph{断言：扩充后的族\omterm{def:b1:vspaces:free}{线性无关}。}
        假设 $\alpha x_{k+1} + \beta f(x_{k+1}) + v = 0$，其中 $v
        \in V_k$ 且 $(\alpha, \beta) \neq (0,0)$。作用 $f$ 得
        $\alpha f(x_{k+1}) - \beta x_{k+1} + f(v) = 0$，其中 $f(v)
        \in V_k$。消去两关系式中的 $f(x_{k+1})$
        （将第一式乘以 $\alpha$，第二式乘以 $-\beta$，
        再相加）：
        \[
        (\alpha^2 + \beta^2)\, x_{k+1} \in V_k ,
        \]
        且 $\alpha^2 + \beta^2 \neq 0$ 迫使 $x_{k+1} \in V_k$：
        矛盾。于是构造得以继续，
        \emph{每次添加两个向量}，直到 $V_k = E$。
        此时最终的族是元素个数为偶数的基，故 $n$ 为偶数。
\end{enumerate}
\end{solution}

\begin{solution}{exo:b1:linmaps:11}
\omterm{def:b1:reals:bounds}{上界}：对每个 $x$，$(u + v)(x) = u(x) + v(x) \in
\operatorname{im} u + \operatorname{im} v$，于是
\[
\operatorname{rk}(u + v)
\leq \dim(\operatorname{im} u + \operatorname{im} v)
\leq \operatorname{rk} u + \operatorname{rk} v
\]
（格拉斯曼，见\cref{thm:b1:findim:grassmann}）。下界：
对二元组 $(u + v, -v)$ 应用\omterm{def:b1:reals:bounds}{上界}，其和为 $u$：
\[
\operatorname{rk} u \leq \operatorname{rk}(u + v) +
\operatorname{rk}(-v) = \operatorname{rk}(u + v) +
\operatorname{rk} v,
\]
由此 $\operatorname{rk} u - \operatorname{rk} v \leq
\operatorname{rk}(u+v)$。交换 $u$ 与 $v$，即得绝对
值形式的不等式。
\end{solution}

\begin{solution}{exo:b1:linmaps:12}
\emph{精确公式。}设 $v'$ 为 $v$ 在
\omterm{def:b1:vspaces:subspace}{子空间} $\operatorname{im} w$ 上的限制，其像为 $v(w(F)) =
\operatorname{im}(v \circ w)$，其核为 $\ker v \cap
\operatorname{im} w$。对 $v'$ 在空间
$\operatorname{im} w$ 上应用秩--零化度定理：
\[
\operatorname{rk} w = \dim\operatorname{im} w
= \operatorname{rk}(v \circ w) + \dim(\ker v \cap
\operatorname{im} w) .
\]

\emph{弗罗贝尼乌斯。}对 $w$ 及
$w \circ u$ 各用一次精确公式：
\[
\operatorname{rk} w - \operatorname{rk}(vw)
= \dim\bigl(\ker v \cap \operatorname{im} w\bigr),
\qquad
\operatorname{rk}(wu) - \operatorname{rk}(vwu)
= \dim\bigl(\ker v \cap \operatorname{im}(wu)\bigr) .
\]
由于 $\operatorname{im}(w \circ u) \subseteq \operatorname{im}
w$，第二个交集包含于第一个，其
维数不更大：
\[
\operatorname{rk}(wu) - \operatorname{rk}(vwu)
\;\leq\; \operatorname{rk} w - \operatorname{rk}(vw) ,
\]
移项即得弗罗贝尼乌斯不等式。取 $w =
\mathrm{id}_F$（秩为 $\dim F$，且 $\operatorname{im}\,
\mathrm{id}_F = F$）：$\operatorname{rk} v + \operatorname{rk} u
\leq \dim F + \operatorname{rk}(vu)$，即西尔维斯特不等式——
在矩阵形式下再次证明，见\cref{exo:b1:matrices:10}。
\end{solution}

\begin{solution}{pb:b1:linmaps:1}
\textbf{1.} $(\mathrm{id} - p)^2 = \mathrm{id} - 2p + p^2 =
\mathrm{id} - p$，故为\omterm{def:b1:linmaps:projection}{投影}。若 $y = x - p(x)$，则 $p(y) =
p(x) - p^2(x) = 0$。反之，$x \in \ker p$ 给出 $x =
(\mathrm{id} - p)(x)$，于是 $\operatorname{im}(\mathrm{id} - p) = \ker
p$。又 $(\mathrm{id} - p)(x) = 0 \iff p(x) = x \iff x \in
\operatorname{im} p$（不动点，
见\cref{thm:b1:linmaps:projchar}）：$\ker(\mathrm{id} - p) =
\operatorname{im} p$。

\textbf{2.} $(\lambda\,\mathrm{id} + \mu p)^2 =
\lambda^2\,\mathrm{id} + (2\lambda\mu + \mu^2)\,p$。二元组
$(\mathrm{id}, p)$ 在 $\mathcal{L}(E)$ 中\omterm{def:b1:vspaces:free}{线性无关}：$p = c\,
\mathrm{id}$ 则 $c^2 = c$，故 $p = 0$ 或 $\mathrm{id}$，
两者均被排除。比较系数，该\omterm{def:b1:logic:map}{映射}为\omterm{def:b1:linmaps:projection}{投影}，当且仅当
$\lambda^2 = \lambda$ 且 $2\lambda\mu + \mu^2 = \mu$。
当 $\lambda = 0$ 时：$\mu \in \{0, 1\}$。当 $\lambda = 1$ 时：$\mu^2 +
\mu = 0$，$\mu \in \{0, -1\}$。平面中恰有四个\omterm{def:b1:linmaps:projection}{投影}，
它们是 $0$，$p$，$\mathrm{id}$，$\mathrm{id} - p$。

\textbf{3.} $(\lambda\,\mathrm{id} + \mu p)(\lambda'\,\mathrm{id}
+ \mu' p) = \lambda\lambda'\,\mathrm{id} + (\lambda\mu' +
\mu\lambda' + \mu\mu')\,p$：该平面在
复合运算下稳定。由于 $p^k = p$ 对每个 $k \geq 1$ 成立，对 $Q =
\sum_k a_k X^k$：
\[
Q(p) = a_0\,\mathrm{id} + \Bigl(\sum_{k \geq 1} a_k\Bigr) p
= Q(0)\,\mathrm{id} + \bigl(Q(1) - Q(0)\bigr)\,p .
\]

\textbf{4.} 在 $\operatorname{im} p$ 上（其中 $p$ 充当
恒等\omterm{def:b1:logic:map}{映射}），$\lambda\,\mathrm{id} + \mu p$ 乘以 $\lambda
+ \mu$；在 $\ker p$ 上，乘以 $\lambda$。由 $E = \operatorname{im} p
\oplus \ker p$，该\omterm{def:b1:logic:map}{映射}\omterm{def:b1:logic:inj}{双射}当且仅当 $\lambda \neq 0$ 且
$\lambda + \mu \neq 0$。解 $\lambda\alpha = 1$，
$\lambda\beta + \mu\alpha + \mu\beta = 0$，在
第 3 题的复合规则中：
\[
(\lambda\,\mathrm{id} + \mu p)^{-1}
= \frac1\lambda\,\mathrm{id} -
\frac{\mu}{\lambda(\lambda + \mu)}\,p ,
\]
其作用为 $1/\lambda$（在 $\ker p$ 上）与 $1/(\lambda +
\mu)$（在 $\operatorname{im} p$ 上），这正是应有的结果。

\textbf{5.} 记 $F = \operatorname{im} p = \operatorname{im}
p'$。对每个 $x$，$p'(x) \in F$，且 $p$ 逐点固定 $F$：
$p(p'(x)) = p'(x)$，即 $p\,p' = p'$；同理 $p'\,p =
p$。当两个\omterm{def:b1:linmaps:projection}{投影}的像相同时，
\emph{先}作用的那个起决定作用：其输出已位于 $F$ 中，在此
外层\omterm{def:b1:linmaps:projection}{投影}充当恒等\omterm{def:b1:logic:map}{映射}，不改变任何东西。

\textbf{6.} $(p + q)^2 = p^2 + pq + qp + q^2 = (p + q) + pq +
qp$，故 $p + q$ 为\omterm{def:b1:linmaps:projection}{投影}迫使 $pq + qp = 0$。在左侧与
$p$ 复合：$pq + pqp = 0$；在右侧与 $p$ 复合：$pqp + qp
= 0$。相减得 $pq = qp$；于是 $pq + qp = 2pq = 0$，而
特征不为 $2$：$pq = qp = 0$。

\textbf{7.} 由 $pq = qp = 0$，同样的展开给出 $(p +
q)^2 = p + q$。像：$\operatorname{im}(p + q) \subseteq
\operatorname{im} p + \operatorname{im} q$ 恒成立。反之，
对 $x \in \operatorname{im} p$：$q(x) = q(p(x)) = 0$，故 $(p +
q)(x) = p(x) = x$，且 $x \in \operatorname{im}(p+q)$；对
$\operatorname{im} q$ 同理。直和性：$x \in \operatorname{im} p
\cap \operatorname{im} q$ 推出 $x = p(x) = p(q(x)) = 0$。
核：若 $p(x) + q(x) = 0$，以 $p$ 作用得 $p(x) +
p(q(x)) = p(x) = 0$，以 $q$ 作用则得 $q(x) = 0$：$\ker(p
+ q) = \ker p \cap \ker q$（反向包含显然成立）。

\textbf{8.} $p - q$ 为\omterm{def:b1:linmaps:projection}{投影}当且仅当 $\mathrm{id} - (p - q) =
(\mathrm{id} - p) + q$ 亦为\omterm{def:b1:linmaps:projection}{投影}（两次应用第 1 题）。由
第 6--7 题应用于\omterm{def:b1:linmaps:projection}{投影} $\mathrm{id} - p$ 与 $q$，此
成立当且仅当 $(\mathrm{id} - p)q = q(\mathrm{id} - p) = 0$，即
当且仅当 $pq = q$ 且 $qp = q$。

\textbf{9.} $pq = q$ 意味着 $p$ 固定每个 $q(x)$，即
$\operatorname{im} q \subseteq \ker(p - \mathrm{id}) =
\operatorname{im} p$。又 $qp = q$ 意味着 $q\bigl((\mathrm{id} -
p)(x)\bigr) = 0$ 对所有 $x$ 成立，即 $q$ 在
$\operatorname{im}(\mathrm{id} - p) = \ker p$ 上为零：$\ker p \subseteq
\ker q$。两步均为等价变换：\omterm{def:b1:logic:order}{序关系} $q \leq p$ 表明
$q$ \omterm{def:b1:linmaps:projection}{投影}到较小的像上，沿较大的核。

\textbf{10.} 展开得 $(\mathrm{id} - p)(\mathrm{id} - q) =
\mathrm{id} - p - q + pq = \mathrm{id} - r$。两个\omterm{def:b1:linmaps:projection}{投影}
$\mathrm{id} - p$ 与 $\mathrm{id} - q$ 可交换，由
\cref{exo:b1:linmaps:6}，其积 $\mathrm{id} - r$ 是
到 $\operatorname{im}(\mathrm{id} - p) \cap
\operatorname{im}(\mathrm{id} - q) = \ker p \cap \ker q$ 上、沿
$\ker(\mathrm{id} - p) + \ker(\mathrm{id} - q) =
\operatorname{im} p + \operatorname{im} q$ 的\omterm{def:b1:linmaps:projection}{投影}。由第 1 题，$r =
\mathrm{id} - (\mathrm{id} - r)$ 则为\omterm{def:b1:linmaps:projection}{投影}，满足
$\operatorname{im} r = \operatorname{im} p + \operatorname{im}
q$ 且 $\ker r = \ker p \cap \ker q$。

\textbf{11.} $p_i$ 是良定义的（由
分解的唯一性），且线性（$x + \lambda y$ 的
分解为各分解之和，同样由唯一性）。对 $x_i
\in F_i$ 时分解即为 $x_i$ 本身，故 $p_i(x_i) = x_i$：
$p_i^2 = p_i$，且 $p_j(x_i) = 0$（$j \neq i$）：$p_i p_j = 0$
（$p_j(x) \in F_j$）。对分量求和，$\sum_i p_i =
\mathrm{id}$。最后 $\operatorname{im} p_i = F_i$，且 $\ker
p_i = \bigoplus_{j \neq i} F_j$。

\textbf{12.} $p_i = p_i \circ \mathrm{id} = p_i\sum_j p_j =
p_i^2 + \sum_{j \neq i} p_i p_j = p_i^2$：每个 $p_i$ 都是
\omterm{def:b1:linmaps:projection}{投影}。任意 $x = \mathrm{id}(x) = \sum_i p_i(x)$ 属于
$\sum_i \operatorname{im} p_i$，即各像之和为 $E$。
直和性：设 $y_1 + \dots + y_k = 0$，其中 $y_i \in
\operatorname{im} p_i$，故 $p_i(y_i) = y_i$。作用 $p_j$：$p_j
(y_i) = p_j p_i (y_i) = 0$（$i \neq j$），于是 $0 = p_j\bigl(\sum
y_i\bigr) = y_j$，对每个 $j$ 成立。因此 $E = \bigoplus_i
\operatorname{im} p_i$。

\textbf{13.} 设 $q = \mathrm{id} - p$，由第 1 题得 $pq =
p - p^2 = 0 = qp$，直接可得：两个\omterm{def:b1:linmaps:projection}{投影}之和为
恒等\omterm{def:b1:logic:map}{映射}，这本身就蕴含二者正交。

\textbf{14.} $p + q = \mathrm{id} - r$，其中 $r$ 为\omterm{def:b1:linmaps:projection}{投影}，且
$(\mathrm{id} - r)$ 为\omterm{def:b1:linmaps:projection}{投影}（第 1 题）：故 $p + q$ 为
\omterm{def:b1:linmaps:projection}{投影}，由第 6 题得 $pq = qp = 0$。由对称性（$q +
r = \mathrm{id} - p$ 与 $p + r = \mathrm{id} - q$），
所有两两乘积皆为零，由第 12 题得出：$E =
\operatorname{im} p \oplus \operatorname{im} q \oplus
\operatorname{im} r$，自动成立。

\textbf{15.} \emph{引理。}由归纳法，结合格拉斯曼公式
（\cref{thm:b1:findim:grassmann}）：
\[
\dim(F_1 + \dots + F_k) \leq \dim(F_1 + \dots + F_{k-1}) + \dim
F_k \leq \dots \leq \sum_i \dim F_i .
\]
若上式取等号，则每一步均取等号，即 $(F_1 + \dots +
F_{j-1}) \cap F_j = \{0\}$ 对每个 $j$ 成立，而关系式 $y_1 +
\dots + y_k = 0$（$y_i \in F_i$）自右端依次推出：$y_k
\in (F_1 + \dots + F_{k-1}) \cap F_k = \{0\}$，于是 $y_{k-1} =
0$ 等等，该和为直和。反之，直和的维数可加
（拼接各基即可）。\emph{应用：}
$x = \sum_i p_i(x)$ 表明 $E = \sum_i \operatorname{im} p_i$，
故 $n \leq \sum_i \operatorname{rk} p_i$；由假设得
等号成立，从而为直和。乘积：固定 $j$ 及 $y \in
\operatorname{im} p_j$。于是 $y = \sum_i p_i(y)$，其中 $p_i(y)
\in \operatorname{im} p_i$，而 $y = y$ 亦是一个
分解（只有第 $j$ 个分量）；由唯一性，$p_i(y)
= 0$（$i \neq j$）。将其用于 $y = p_j(x)$：$p_i p_j = 0$。

\textbf{16.} 若 $u^k(x) = 0$，则 $u^{k+1}(x) = u(0) = 0$。又有
$\operatorname{im} u^{k+1} = u^k\bigl(u(E)\bigr) \subseteq
u^k(E) = \operatorname{im} u^k$.

\textbf{17.} 设 $\ker u^{r} = \ker u^{r+1}$，取 $x \in
\ker u^{r+2}$，则 $u(x) \in \ker u^{r+1} = \ker u^{r}$，于是
$u^{r+1}(x) = 0$，即 $x \in \ker u^{r+1}$。由第 16 题，
$\ker u^{r+1} = \ker u^{r+2}$，由归纳法，此后
各核均与 $\ker u^{r}$ 相同。对于像，由秩--零化度定理
得 $\dim\operatorname{im} u^k = n - \dim\ker u^k$，故
像的维数恰与核的维数同时稳定，
结合第 16 题中的包含关系，维数相等即
\omterm{def:b1:vspaces:subspace}{子空间}相等（由\cref{thm:b1:findim:subspaces}可知）。

\textbf{18.} 数列 $\bigl(\dim\ker u^k\bigr)_k$ 是
单调不减，取值于 $\intint{0}{n}$；它不可能
严格增加 $n + 1$ 次，故存在 $\dim\ker u^{r} = \dim\ker
u^{r+1}$，其中 $r \leq n$，从而 $\ker u^{r} = \ker u^{r+1}$
（包含关系加维数相等）。取 $r$ 为最小者，由第 17 题
可知，自 $r$ 起一切（包括像）均保持不变。

\textbf{19.} 先看交：取 $x \in \ker u^{r} \cap
\operatorname{im} u^{r}$，记 $x = u^{r}(y)$，且 $u^{r}(x) =
0$。于是 $u^{2r}(y) = 0$，而 $\ker u^{2r} = \ker u^{r}$
（由第 17 题），故 $x = u^{r}(y) = 0$。维数方面，由秩--零化度定理
对 $u^{r}$ 得 $\dim\ker u^{r} + \dim\operatorname{im} u^{r}
= n$；交为零，由格拉斯曼公式知和
是维数为 $n$ 的\omterm{def:b1:vspaces:subspace}{子空间}：$E = \ker u^{r} \oplus
\operatorname{im} u^{r}$。

\textbf{20.} 稳定性：$u^{r}(u(x)) = u(u^{r}(x)) = 0$，对 $x
\in \ker u^{r}$；又 $u(u^{r}(y)) = u^{r}(u(y)) \in
\operatorname{im} u^{r}$。在 $N = \ker u^{r}$ 上，$(u|_N)^{r} = 0$
由 $N$ 的定义，即幂零。在 $I = \operatorname{im}
u^{r}$: $\ker(u|_I) = \ker u \cap I \subseteq \ker u^{r} \cap I
= \{0\}$，故 $u|_I$ 是
\omterm{def:b1:findim:def}{有限维}空间 $I$ 上的单自\omterm{def:b1:structures:morphism}{同态}，从而为\omterm{def:b1:logic:inj}{双射}
（由\cref{cor:b1:linmaps:samedim}可知）。

\textbf{21.} 设 $x = a + b$，其中 $a \in N$，$b \in I$。则
$u(x) = u(a) + u(b)$，且 $u(a) \in N$，$u(b) \in I$
（由第 20 题），这\emph{正是} $u(x)$ 的分解，故
$\pi(u(x)) = u(a) = u(\pi(x))$: $\pi u = u\pi$.

\textbf{22.} $u^2(x,y,z) = u(y, 0, z) = (0, 0, z)$，且
$u^3(x,y,z) = u(0,0,z) = (0,0,z) = u^2(x,y,z)$。核分别为 $\ker u
= \{y = z = 0\} = \operatorname{Vect}(e_1)$, $\ker u^2 = \{z =
0\} = \operatorname{Vect}(e_1, e_2)$, $\ker u^3 = \ker u^2$:
在 $r = 2$ 处稳定。像：$\operatorname{im} u =
\operatorname{Vect}(e_1, e_3)$, $\operatorname{im} u^2 =
\operatorname{Vect}(e_3)$。费廷分解：$\R^3 =
\operatorname{Vect}(e_1, e_2) \oplus \operatorname{Vect}(e_3)$,
且 $\pi(x, y, z) = (x, y, 0)$。验证：$\pi u(x,y,z) = \pi(y, 0,
z) = (y, 0, 0)$，而 $u\pi(x,y,z) = u(x, y, 0) = (y, 0, 0)$：
相等。在第一个因子上 $u(x, y, 0) = (y, 0, 0)$，其
平方为 $0$，故幂零；在第二个因子上 $u(0,0,z) = (0,0,z)$：
恒等\omterm{def:b1:logic:map}{映射}，因而\omterm{def:b1:logic:inj}{双射}。

\textbf{23.} 若 $u^m = 0$，则 $\ker u^m = E$；由于核
自 $r$ 起保持不变，故 $\ker u^{r} = \ker u^{\max(m, r)} = E$。
反之，$\ker u^{r} = E$ 意味着 $u^{r} = 0$。又 $\ker u^{r} =
E \iff$ 等价于费廷\omterm{def:b1:linmaps:projection}{投影}是以 $E$ 为像、以 $\{0\}$ 为核的\omterm{def:b1:linmaps:projection}{投影}，即\
$\pi = \mathrm{id}$。最后由第 18 题得 $r \leq n$，于是：
每个幂零自\omterm{def:b1:structures:morphism}{同态}都满足 $u^{n} = 0$——亦即
幂零指数不超过空间维数。

\textbf{24.} 设 $m$ 为 $u|_A$ 的一个幂零指数：$A
\subseteq \ker u^{m} \subseteq \ker u^{\max(m,r)} = \ker u^{r}$.
由于 $u|_B$ 是\omterm{def:b1:logic:inj}{双射}，$B = u(B) = u^{k}(B) \subseteq
\operatorname{im} u^{k}$ 对每个 $k$ 成立，特别地 $B
\subseteq \operatorname{im} u^{r}$。于是
\[
n = \dim A + \dim B \leq \dim\ker u^{r} +
\dim\operatorname{im} u^{r} = n :
\]
两个包含关系的维数都取等号，从而
\omterm{def:b1:vspaces:subspace}{子空间}亦相等：$A = \ker u^{r}$，$B = \operatorname{im} u^{r}$。

\textbf{25.} (i) 第三部分是一部“词典”：分解 $E = F_1
\oplus \dots \oplus F_k$ 恰好对应于
满足 $\sum p_i = \mathrm{id}$ 与 $p_i p_j = 0$ 的\omterm{def:b1:linmaps:projection}{投影}族，
$F_i$ 为各\omterm{def:b1:linmaps:projection}{投影}的像。(ii) 当 $k = 3$ 时，补
$\mathrm{id} - p_i$ 本身即为\omterm{def:b1:linmaps:projection}{投影}，这恰好
完成了论证，无需额外假设；对一般的 $k$，需要
$\sum_i \operatorname{rk} p_i \leq n$，这一不等式可
由本科二年级的迹直接得到（对\omterm{def:b1:linmaps:projection}{投影}有 $\operatorname{tr} p = \operatorname{rk}
p$，且迹相加得 $\operatorname{tr}
\mathrm{id} = n$）。(iii) \cref{exo:b1:linmaps:7} 是费廷
引理在已稳定情形 $r \leq 1$ 下的特例；一般情形下
让核链与像链稳定，至多需要 $n$
步。(iv) 在本科二年级分册中，将其应用于 $u - \lambda\,
\mathrm{id}$，幂零因子变为
$\lambda$ 处的广义特征\omterm{def:b1:vspaces:subspace}{子空间}，第三部分的\omterm{def:b1:linmaps:projection}{投影}则变为
约化理论中的谱\omterm{def:b1:linmaps:projection}{投影}。第四部分的定理
即\emph{费廷引理}。
\end{solution}
